%% =====================================================================
\chapter{The Zero-Trust Layer}\label{ch:zerotrust}

``Zero trust'' means: never assume a request is legitimate because of an
earlier check. MediChain applies four checks to \emph{every} protected request,
in the middleware in \file{src\codesl{}middleware\codesl{}auth\codedot{}js} and the policy engine in
\file{src\codesl{}services\codesl{}trustService\codedot{}js}, and records every decision.

\begin{enumerate}
  \item \textbf{Authentication}: is the token genuine, unexpired and not revoked?
  \item \textbf{Continuous verification}: is the account still active, and what
    is its current state?
  \item \textbf{Trust score}: how much does this particular request look like
    its legitimate owner?
  \item \textbf{Authorisation}: does the role, approval status, trust score and
    (for patient data) consent allow this action?
\end{enumerate}

\section{\code{authenticate}}

\begin{lstlisting}[language=JavaScript,caption={authenticate (\file{src/middleware/auth.js})}]
async function authenticate(req, res, next) {
  const header = req.headers.authorization || '';
  if (!header.startsWith('Bearer ') || !header.slice(7).trim()) {
    return deny(req, res, 401, 'missing_token', 'Access denied. No token provided. Please log in.');
  }

  let decoded;
  try {
    decoded = await tokenService.verifyAccessToken(header.slice(7).trim());
  } catch (error) {
    if (error.name === 'TokenExpiredError') return deny(req, res, 401, 'token_expired', 'Token expired. Please log in again.');
    if (error.message === 'TOKEN_REVOKED') return deny(req, res, 401, 'token_revoked', 'Token has been revoked. Please log in again.');
    if (error.message === 'REDIS_UNAVAILABLE') return deny(req, res, 503, 'redis_unavailable', 'Security service temporarily unavailable. Please try again.');
    return deny(req, res, 401, 'invalid_token', 'Invalid token. Please log in again.');
  }

  let sessionState;
  try {
    sessionState = await loadSessionState(decoded.userId);
  } catch (error) {
    logger.error(`Session state unavailable: ${error.message}`);
    return deny(req, res, 503, 'session_cache_unavailable', 'Security service temporarily unavailable. Please try again.', { userId: decoded.userId });
  }

  if (!sessionState || !sessionState.isActive) {
    return deny(req, res, 401, 'account_deactivated', 'Account deactivated. Contact the hospital administrator.', { userId: decoded.userId });
  }

  req.user = decoded; // { userId, email, role, jti, iat, exp }
  req.sessionState = sessionState;
  req.trustScore = trustService.calculateTrustScore(req, decoded, sessionState);

  if (req.trustScore.score < req.trustScore.minScore) {
    return deny(req, res, 403, 'low_trust_score', 'Access denied. Additional verification required: please log in again.', {
      userId: decoded.userId, trustScore: req.trustScore.score, body: { trustScore: req.trustScore.score },
    });
  }

  await auditService.recordAccessDecision({ req, userId: decoded.userId, decision: 'allow', reason: 'authenticated', trustScore: req.trustScore.score });
  return next();
}
\end{lstlisting}

\begin{howto}
Every failure path calls \code{deny}, which writes the decision to the ledger's
access chain before answering, so failed attempts are as visible as successful
ones. The reason code (\code{token\_expired}, \code{token\_revoked}, \ldots) is
returned to the client in upper case; the frontend uses \code{401} to try a
silent refresh and \code{503} to show ``security service unavailable''.
\end{howto}

\subsection{Continuous verification with a session cache}

\begin{lstlisting}[language=JavaScript,caption={loadSessionState (\file{src/middleware/auth.js})}]
async function loadSessionState(userId) {
  const cached = await redisClient.get(sessionKey(userId));
  if (cached) return JSON.parse(cached);

  const user = await User.findById(userId).select('isActive isVerified failedLoginAttempts approvalStatus').lean();
  if (!user) return null;
  const state = {
    isActive: user.isActive,
    isVerified: user.isVerified,
    failedLoginAttempts: user.failedLoginAttempts || 0,
    approvalStatus: user.approvalStatus,
  };
  await redisClient.setEx(sessionKey(userId), SESSION_CACHE_SECONDS, JSON.stringify(state));
  return state;
}
\end{lstlisting}

A JWT says who the user \emph{was} when it was issued. If an admin deactivates
the account a minute later, the token is still cryptographically valid for 14
more minutes. Loading the account state on each request closes that gap. To
avoid a MongoDB query per request, the state is cached in Redis for 5 minutes,
and every code path that changes it (deactivation, approval, failed login,
password change) calls \code{invalidateSession}, so changes take effect on the
very next request.

\section{The trust score}\label{sec:trust}

\begin{lstlisting}[language=JavaScript,caption={calculateTrustScore (\file{src/services/trustService.js})}]
  calculateTrustScore(req, tokenPayload, sessionState = {}) {
    let score = 100;
    const reasons = [];
    const failedAttempts = Number(sessionState.failedLoginAttempts || 0);
    const sessionAgeSeconds = tokenPayload?.iat
      ? Math.max(0, Math.floor(Date.now() / 1000) - tokenPayload.iat)
      : null;
    const userAgent = req.get('user-agent') || '';

    if (!sessionState.isVerified) {
      score -= 15;
      reasons.push('user_not_verified');
    }
    if (failedAttempts > 0) {
      score -= Math.min(30, failedAttempts * 6);
      reasons.push('recent_failed_attempts');
    }
    if (!userAgent) {
      score -= 10;
      reasons.push('missing_user_agent');
    }
    if (sessionAgeSeconds !== null && sessionAgeSeconds > 60 * 60) {
      score -= 10;
      reasons.push('older_session');
    }
    if (sessionAgeSeconds !== null && sessionAgeSeconds > 6 * 60 * 60) {
      score -= 15;
      reasons.push('long_lived_session');
    }
    // ... returns { score, reasons, minScore, highSensitivityMinScore, ... }
\end{lstlisting}

\begin{table}[H]
\centering\small
\caption{Trust score signals}\label{tab:trust}
\begin{tabularx}{\textwidth}{@{}P{4.2cm}P{1.6cm}Y@{}}
\toprule
\textbf{Signal} & \textbf{Penalty} & \textbf{Why it matters} \\
\midrule
Email never verified by OTP & $-15$ & Weaker proof of identity \\
Recent failed logins & $-6$ each, max $-30$ & Someone may be guessing this account's password \\
No User-Agent header & $-10$ & Typical of scripts rather than browsers \\
Session older than 1 hour & $-10$ & Old credentials are more likely to have leaked \\
Session older than 6 hours & $-15$ more & As above \\
\bottomrule
\end{tabularx}
\end{table}

Every request needs at least \code{MIN\_TRUST\_SCORE} (50). Medical records and
files need \code{HIGH\_SENSITIVITY\_TRUST\_SCORE} (70). Example: a user with two
recent failed logins ($-12$) whose access token was issued 2 hours ago via
refresh ($-10$) scores 78: allowed everywhere. With five failures ($-30$) and a
7-hour-old token ($-25$) the score is 45: every request is refused with
``please log in again'', and a fresh login (which resets the failures) restores
access. The profile page shows each user their current score and the reasons.

\section{The policy engine}\label{sec:policy}

\begin{lstlisting}[language=JavaScript,caption={evaluatePolicy (\file{src/services/trustService.js})}]
  evaluatePolicy(req, allowedRoles, options = {}) {
    const role = req.user?.role;
    const trustScore = req.trustScore?.score ?? 0;
    const minTrustScore = options.minTrustScore ?? config.security.minTrustScore;
    const breakGlassReason = (req.get('x-break-glass-reason') || '').trim();
    const breakGlass = Boolean(options.allowBreakGlass && breakGlassReason.length >= 10 && BREAK_GLASS_ROLES.includes(role));

    if (breakGlass) {
      return { allowed: true, reason: 'break_glass_allowed', breakGlass: true, breakGlassReason };
    }
    if (!allowedRoles.includes(role)) {
      return { allowed: false, status: 403, reason: 'role_denied', message: `Access denied. This resource requires the ${allowedRoles.join(' or ')} role.` };
    }
    if (ROLES_NEEDING_APPROVAL.includes(role) && !options.allowUnapproved && req.sessionState?.approvalStatus !== 'approved') {
      return { allowed: false, status: 403, reason: 'approval_pending', message: 'Your account is awaiting administrator verification.' };
    }
    if (options.requireVerified && !req.sessionState?.isVerified) {
      return { allowed: false, status: 403, reason: 'verification_required', message: 'Access denied. Verified account required.' };
    }
    if (trustScore < minTrustScore) {
      return { allowed: false, status: 403, reason: 'low_trust_score', message: `Access denied. This resource needs a trust score of ${minTrustScore} (yours is ${trustScore}). Log in again to refresh your session.` };
    }
    if (options.timeWindow && !this.isWithinTimeWindow(new Date().getHours(), options.timeWindow)) {
      return { allowed: false, status: 403, reason: 'outside_time_window', message: 'Access denied outside the approved access window.' };
    }
    return { allowed: true, reason: 'policy_allowed', breakGlass: false };
  }
\end{lstlisting}

\begin{howto}
Routes declare their policy declaratively:
\begin{lstlisting}[language=JavaScript]
const readPolicy = requireRole('patient', 'doctor', {
  minTrustScore: config.security.highSensitivityTrustScore, allowBreakGlass: true,
});
\end{lstlisting}
\code{requireRole} (in \file{auth\codedot{}js}) calls \code{evaluatePolicy} and either
continues or calls \code{deny}. The checks run in a fixed order and the first
failure wins: role (RBAC), then approval, verification, trust score and time
window (attribute-based conditions, ABAC). The result is stored in
\code{req.policy} so the handler knows whether break-glass was used.
\end{howto}

\section{Break-glass emergency access}

An unconscious patient in the emergency room cannot approve a consent request.
Break-glass lets a doctor (or admin) open the record anyway, under three
conditions: the route must allow it (\code{allowBreakGlass}), the user must be a
doctor or admin, and the request must carry a reason of at least 10 characters
in the \code{x-break-glass-reason} header. Break-glass replaces the
\emph{consent} check only; the user must still be authenticated.

What happens next is in \code{consentService.assertCanAccessPatient}
(Chapter~\ref{ch:records}): the access is written to the records chain as
\code{BREAK\_GLASS\_ACCESS} with the reason, the patient gets a ``danger''
notification quoting the reason, and every admin is notified to review it.
Emergency access is therefore possible, but never silent.

\section{The audit trail: \file{src\codesl{}services\codesl{}auditService\codedot{}js}}

\begin{lstlisting}[language=JavaScript,caption={recordAccessDecision (\file{src/services/auditService.js})}]
async function recordAccessDecision({ req, userId, decision, reason, trustScore, breakGlassReason, resource }) {
  try {
    await ledgerService.append({
      chain: 'access',
      entityType: 'access_decision',
      entityId: resource || `${req.method} ${req.originalUrl || req.path}`,
      action: decision === 'allow' ? 'ALLOW' : 'DENY',
      actorId: userId || req.user?.userId || null,
      actorRole: req.user?.role || null,
      meta: {
        reason,
        method: req.method,
        path: (req.originalUrl || req.path || '').split('?')[0],
        ip: req.ip,
        userAgentHash: sha256(req.get('user-agent') || ''),
        trustScore: trustScore ?? req.trustScore?.score ?? null,
        breakGlassReason: breakGlassReason || undefined,
      },
    });
  } catch (error) {
    logger.error(`Audit write failed: ${error.message}`);
  }
}
\end{lstlisting}

Each decision becomes one block of the \code{access} chain, so the audit log is
itself tamper-evident. The query string is dropped from the path and the
User-Agent is stored only as a hash, to avoid putting personal data in a log
that can never be edited. Audit failures are logged but do not fail the request:
losing one audit line is better than locking every user out. The admin
dashboard's \emph{Access Control} page reads this chain through
\code{GET \codesl{}api\codesl{}admin\codesl{}audit}.
